% Ad Familiares 7.20 (ad-familiares-07-20) --- English translation
% Translated by Claude (Anthropic) from the Latin (Perseus).
% Style guide: STYLE.md.

\ciceroLetterOpener{CICERO TREBATIO S.}

\ciceroSection{1} Velia was the dearer to me because I felt how dear you are to it. But why do I say ``you,'' whom does not love? Your Rufio, by Hercules, was missed as much as if he were one of ourselves. Still, I do not fault you for hauling him off to your building project. For though Velia is no cheaper than the Lupercal, I still prefer where you are to all this here. If you take my advice --- which you usually do --- you will hold on to these paternal properties (the Velians were worried about something or other), and you will not leave the Halaes, that noble river, nor desert the Papirian house --- though that house, to be sure, has a lotus by which even strangers are commonly held captive; and yet if you cut it down, you will both see further and have shown foresight.

\ciceroSection{2} But this above all looks well-judged, particularly in times like these: to have a place of refuge --- first, a town where you are held dear, then your own house and your own lands, and those in a remote, healthy, agreeable spot. And I think this concerns me too, my Trebatius. But farewell. Look after my affairs, and, with the gods' help, expect me back before midwinter.

\ciceroSection{3} I have carried off from Sex. Fadius, Nico's pupil, a book of Nico's \textit{On Overeating} \textit{[Greek: Nikōnos peri polyphagias]}. What a delightful doctor --- and how teachable I prove in this discipline! But our friend Bassus kept this book hidden from me --- though apparently not from you. The wind is rising. Take care of yourself. The thirteenth day before the Kalends of Sextilis. Velia.
